\section{Introduction}

\begin{frame}{文献简介}
  \begin{itemize}
    \item 标题：Digital Collateral
    \item 作者：Paul Gertler (UC Berkeley), Brett Green (Wash U), Catherine Wolfram (MIT Sloan)
    \item 期刊：\textit{The Quarterly Journal of Economics} (Top 5)
    \item 年份：2024
    \item \cite{gertler2024DigitalCollateral}
  \end{itemize}
\end{frame}

\begin{frame}[allowframebreaks]{Motivation}
  \begin{block}{核心问题}
    \textbf{抵押品（collateral）是信贷市场的基石——但在发展中国家，传统抵押品为什么不好用？}
  \end{block}

 \medskip

  \begin{block}{供给侧的障碍}
  \begin{itemize}
    \item 弱法治环境下产权难以界定和执行
    \item 物理回收成本极高——偏远地区的借款人很难被找到，找到之后资产可能已经损坏或转移
    
    \begin{itemize}
      \item 美国超过80\%的家庭债务有物理资产担保
      \item 我们的初步调查显示，只有 30\% 的家庭贷款以实际资产（如房屋、土地、牲畜或车辆）为抵押。 
      \item 乌干达“透明定价倡议”的数据显示，在 2011-2014 年期间，仅 16\% 的贷款产品明确提及使用抵押品。 
      \item Elum 和 Kanu (2018) 的调查发现，尼日利亚 47\% 的农民认为抵押品是获得正式信贷的主要障碍。
    \end{itemize}
  \end{itemize}
  \end{block}

 \medskip

  \begin{block}{需求侧的障碍}
  \begin{itemize}
    \item 许多中低收入国家家庭的主要收入来源是自营职业,这比正规部门的工资更容易受到频繁冲击。
    \item 这些缺乏储蓄的家庭更可能因非策略性原因违约,  并可能选择避免资产被收回的风险。
  \end{itemize}
\end{block}
  \medskip

  \begin{block}{传统抵押品在LMIC的现实困境}
  回收成本高（$\kappa$ 低）、借款人面临永久丧失资产的风险 $\to$ 担保债务的普及程度极低
  \end{block}
\end{frame}

\begin{frame}{数字抵押品 --- 一种新的解法}
  \begin{block}{核心创新}
  \textbf{lockout technology（锁定技术）}
  

  \begin{itemize}
    \item 贷款人可以\textbf{远程禁用}资产的服务流（flow of services），而无需物理回收
    \item 禁用成本极低、可逆，即借款人补缴后立即恢复使用
    \item 典型应用场景：Pay-As-You-Go（PAYGO）模式：小额首付 → 取得资产 → 通过移动支付分期还款 → 逾期即锁定 → 补缴即解锁
  \end{itemize}

  \end{block}

  \begin{block}{传统抵押品的双重功能}
    \begin{itemize}
        \item 贷款方能回收有价值资产，从而规避违约风险
        \item 借款家庭将面临价值损失，这既激励其履行还款义务，也促使预期违约者主动拒绝贷款
    \end{itemize}
  \end{block}
\end{frame}

\begin{frame}{数字抵押品的应用谱系——为什么这和"数字金融"有关}
  这不仅是一个发展经济学的故事，更是数字技术重塑金融合约基础设施的典型案例。

  \begin{block}{已有的应用场景}
    \begin{itemize}
      \item 太阳能家庭系统 (SHS)：Fenix International（乌干达），本文实验对象
      \item 公用事业：TELMEX（墨西哥电信）用网络服务作为数字抵押品为电脑设备贷款提供担保
      \item 汽车贷款：美国已在次级汽车贷款领域应用了类似技术——超过两百万辆汽车安装了"启动阻断装置"，当借款人未按时还款时，贷款方可远程禁用车辆启动功能。
    \end{itemize}
  \end{block}
  
\medskip

  \begin{block}{中国场景的类比}
    \begin{itemize}
      \item 共享充电宝/共享单车的信用免押机制
      \item 手机分期的锁机技术（如趣分期早期模式，尽管该模式引发了较大的舆论争议） %TODO 验证这一说法
      \item 芝麻信用的"信用即服务"逻辑
      \item 社会信用体系下的失信联合惩戒
    \end{itemize}
  \end{block}
\end{frame}

\begin{frame}[allowframebreaks]{问题}
    \begin{exampleblock}{问题1：}
    \textbf{数字抵押与传统抵押的关键区别}
    \end{exampleblock}
  \begin{table}
  \centering
  \caption{传统抵押品 vs. 数字抵押品}
  \label{tab:traditional_vs_digital_collateral}  % Add a label for referencing
  {\begin{tabular}{lll}
    \toprule
    \textbf{维度} & \textbf{传统抵押品} & \textbf{数字抵押品} \\
    \midrule
    回收方式 & 物理没收 & 远程锁定 \\
    回收成本 & 高（交通、法律、处置） & 极低（远程操作） \\
    可逆性 & 不可逆（永久丧失） & 可逆（补缴即恢复） \\
    \rowcolor{yellow!30}
    贷款人回收价值 & 有 & 几乎没有 \\
    对借款人的激励 & 强（怕失去资产） & 强（怕失去服务流） \\
    \bottomrule
  \end{tabular}}
\end{table}


  
   \begin{exampleblock}{问题2：}
      \textbf{乌干达的当地百姓在使用数字抵押品获得抵押贷款后，这笔资金被用在了什么地方？}
    \end{exampleblock}

    \medskip

   \begin{exampleblock}{问题3：}
      \textbf{这样的数字抵押品“好”还是“不好”？}
      \begin{itemize}
        \item 是否有助于还款义务的履行？
        \item 是否会恶化家庭的资产负债表？
        \item 数字担保贷款是否构成一个可持续的（profitable）均衡合约？ %TODO 搞懂什么是可持续的（profitable）均衡合约
        \item 是否改善了福利水平？
      \end{itemize}
      
      
  \end{exampleblock}
\end{frame}